% =====================================================================
% Problem statement -- carries over the content of the approved topic
% proposal. Please do not change the structure.
% =====================================================================
\vorspannkapitel{Problemstellung}
\abschnittsautor{\alleautoren}

\vorspannabschnitt{Ausgangslage}

Bestehende Erkennungssysteme identifizieren bekannte Angriffe häufig anhand fester Regeln. Unbekannte oder veränderte Angriffsmuster können dadurch unentdeckt bleiben. Für eine wirksame Reaktion müssen Netzwerkereignisse daher erfasst, datenbasiert bewertet und in nachvollziehbare, kontrolliert umsetzbare Maßnahmen überführt werden.

\vorspannabschnitt{Untersuchungsanliegen der individuellen Themenstellungen}

\vorspannunterpunkt{1) Teilbereich A}
Untersucht wird, wie eine isolierte Testumgebung mit überwachten Zielsystemen aufgebaut und Netzwerkverkehr zuverlässig sowie technologieoffen erfasst werden kann. Kontrollierte Angriffsszenarien werden geplant, reproduzierbar ausgeführt und zur Bewertung regel- oder signaturbasierter Erkennungsverfahren eingesetzt. Untersucht werden Datenqualität, Erkennungsleistung, Fehlalarme und Ressourcenbedarf. Die aufbereiteten Ereignisse und Alarme werden über definierte Schnittstellen bereitgestellt.

\vorspannunterpunkt{2) Teilbereich B}
Untersucht wird, welche aus Netzwerk- und Sicherheitsereignissen abgeleiteten Merkmale zur Unterscheidung zwischen normalem und auffälligem Verhalten geeignet sind. Mehrere Verfahren zur Anomalieerkennung werden mit Normalverkehr und definierten Angriffsszenarien trainiert beziehungsweise parametrisiert und verglichen. Regelbasierte Alarme und ML-Ergebnisse werden korreliert, zu Vorfällen zusammengeführt und über einen nachvollziehbaren Risiko-Score bewertet. Bewertet werden Erkennungsrate, Fehlalarme, Laufzeit und Ressourcenbedarf.

\vorspannunterpunkt{3) Teilbereich C}
Untersucht wird, wie bewertete Sicherheitsvorfälle mithilfe eines Sprachmodells verständlich erklärt, priorisiert und mit technischen sowie organisatorischen Maßnahmen verknüpft werden können. Das Modell soll strukturierte Regelvorschläge erzeugen. Ein Dashboard stellt Ereignisse, Risiko-Scores, Begründungen, Maßnahmen und Regelvorschläge dar. Prüfung, Test, Freigabe und Aktivierung erfolgen über kontrollierte Prozesse, damit keine ungeprüften Änderungen übernommen werden.

\vorspannabschnitt{Zielsetzung}

Ziel ist ein modulares System, das Netzwerkereignisse erfasst, regel- und datenbasiert bewertet, verständlich interpretiert und kontrollierte Gegenmaßnahmen vorschlägt. Die Komponenten werden technologieoffen entwickelt, über definierte Schnittstellen verbunden und anhand reproduzierbarer Szenarien evaluiert.
\vorspannabschnitt{Geplantes Ergebnis der individuellen Themenstellungen}

\vorspannunterpunkt{1) Teilbereich A}
Es entsteht eine isolierte, reproduzierbare Testumgebung mit überwachten Zielsystemen, normalem Testverkehr und dokumentierten Angriffsszenarien. Eine modulare Komponente erfasst und vereinheitlicht Netzwerkdaten. Mindestens ein regelbasiertes Erkennungsverfahren erzeugt Alarme, die zusammen mit den aufbereiteten Ereignissen über dokumentierte Schnittstellen bereitgestellt und hinsichtlich Erkennungsqualität, Fehlalarmen und Systembelastung evaluiert werden.

\vorspannunterpunkt{2) Teilbereich B}
Es entsteht eine Analysekomponente zur ML-gestützten Anomalieerkennung. Sie verarbeitet aufbereitete Netzwerkdaten, vergleicht mehrere Verfahren und weist auffällige Merkmale aus. Regelbasierte Alarme und ML-Ergebnisse werden zu Sicherheitsvorfällen korreliert und mit einem gemeinsamen Risiko-Score bewertet. Modelle, Datenpipeline, Evaluation und eine dokumentierte Schnittstelle zur Bereitstellung der Vorfälle werden versioniert und getestet.

\vorspannunterpunkt{3) Teilbereich C}
Es entsteht eine Komponente zur sprachmodellgestützten Interpretation bewerteter Sicherheitsvorfälle sowie ein Dashboard zur Darstellung und Bearbeitung. Das System erzeugt verständliche Erklärungen, priorisierte Maßnahmen und strukturierte Regelvorschläge. Ein kontrollierter Ablauf prüft Vorschläge, testet Änderungen in einer isolierten Umgebung und fordert vor der Aktivierung eine Freigabe an. Vorschläge, Prüfungen und Entscheidungen werden nachvollziehbar protokolliert.

\vorspannabschnitt{Meilensteine}

% Deliberately not a table float: the milestones belong directly under the
% heading, not somewhere else in the document.
\begin{tabular}{@{}p{0.68\linewidth}p{0.24\linewidth}@{}}
    \toprule
    \bfseries Meilenstein & \bfseries Datum \\
    \midrule
    Antragsabgabe und Projektstart & 18.09.2026 \\
    Anforderungen, Gesamtarchitektur, Datenformate und Schnittstellen festgelegt & 09.10.2026 \\
    Isolierte Testumgebung, Zielsysteme, Datenerfassung und Angriffsszenarien konzipiert & 06.11.2026 \\
    Erste Prototypen für Erkennung, ML-Analyse, Ereignisinterpretation und Dashboard umgesetzt & 04.12.2026 \\
    Trainings- und Testdaten aufbereitet sowie Kernverfahren implementiert & 15.01.2026 \\
    Ereigniskorrelation, Risiko-Score und durchgängiger Systemprototyp fertiggestellt & 12.02.2026 \\
    Modellvergleich, Regelvalidierung, Dashboard und Gesamtevaluation abgeschlossen & 12.03.2027 \\
    Implementierung abgeschlossen und vollständiger Dokumentationsentwurf erstellt & 09.04.2027 \\
    Abschließende Systemtests, Fehlerbehebung und formale Kontrolle abgeschlossen & 23.04.2027 \\
    Abgabe der Diplomarbeit  & 03.05.2027 \\
  \bottomrule
\end{tabular}
